
\thispagestyle{empty}
\begin{center}
  \vspace*{14mm}
  {\Large\MakeUppercase{Contents}}
\end{center}
\vspace{16mm}

\begingroup
\small
\GuideChapter{\textbf{XIX. Reductive Groups---Generalities}, by
  M.~Demazure}{1}
\GuideSection{1. Review of groups over an algebraically closed field}{1}
\GuideSection{2. Reductive group schemes. Definitions and first
  properties}{8}
\GuideSection{3. Roots and root systems of reductive group schemes}{12}
\GuideSection{4. Roots and vector group schemes}{15}
\GuideSection{5. An instructive example}{20}
\GuideSection{6. Local existence of maximal tori. The Weyl group}{23}
\GuideSection{Bibliography}{25}

\GuideChapter{\textbf{XX. Reductive Groups of Semisimple Rank \(1\)}, by
  M.~Demazure}{27}
\GuideSection{1. Elementary systems. The groups \(U_{\alpha}\) and
  \(U_{-\alpha}\)}{27}
\GuideSection{2. Structure of elementary systems}{35}
\GuideSection{3. The Weyl group}{44}
\GuideSection{4. The isomorphism theorem}{51}
\GuideSection{5. Examples of elementary systems; applications}{53}
\GuideSection{6. Generators and relations for an elementary system}{58}

\GuideChapter{\textbf{XXI. Root Data}, by M.~Demazure}{63}
\GuideSection{1. Generalities}{63}
\GuideSection{2. Relations between two roots}{67}
\GuideSection{3. Simple roots, positive roots}{71}
\GuideSection{4. Reduced root data of semisimple rank \(2\)}{84}
\GuideSection{5. The Weyl group: generators and relations}{86}
\GuideSection{6. Morphisms of root data}{90}
\GuideSection{7. Structure}{101}
\GuideSection{Bibliography}{108}

\GuideChapterSplit{XXII. Reductive Groups: Splittings, Subgroups,
  Quotient Groups}{by M.~Demazure}{109}
\GuideSection{1. Roots and coroots. Split groups and root data}{109}
\endgroup

\clearpage
\begingroup
\small
\GuideSection{2. Existence of a splitting. Type of a reductive group}{114}
\GuideSection{3. The Weyl group}{116}
\GuideSection{4. Homomorphisms of split groups}{118}
\GuideSection{5. Subgroups of type \textup{(R)}}{128}
\GuideSection{6. The derived group}{168}
\GuideSection{Bibliography}{175}

\GuideChapter{\textbf{XXIII. Reductive Groups: Uniqueness of Pinned
  Groups}, by M.~Demazure}{177}
\GuideSection{1. Pinnings}{177}
\GuideSection{2. Generators and relations for a pinned group}{182}
\GuideSection{3. Groups of semisimple rank \(2\)}{190}
\GuideSection{4. Uniqueness of pinned groups: fundamental theorem}{202}
\GuideSection{5. Corollaries of the fundamental theorem}{206}
\GuideSection{6. Chevalley systems}{209}
\GuideSection{Bibliography}{213}

\GuideChapter{\textbf{XXIV. Automorphisms of Reductive Groups}, by
  M.~Demazure}{215}
\GuideSection{1. Scheme of automorphisms of a reductive group}{216}
\GuideSection{2. Automorphisms and subgroups}{224}
\GuideSection{3. Dynkin scheme of a reductive group. Quasi-split
  groups}{228}
\GuideSection{4. Isotriviality of reductive groups and of principal
  bundles under reductive groups}{237}
\GuideSection{5. Canonical decomposition of an adjoint or simply
  connected group}{243}
\GuideSection{6. Automorphisms of Borel subgroups of reductive
  groups}{248}
\GuideSection{7. Representability of the functors
  \(\operatorname{Hom}_{S\text{-grp.}}(G,H)\), for \(G\) reductive}{250}
\GuideSection{8. Appendix: Cohomology of a smooth group over a henselian
  ring. Cohomology and the functor \(\prod\)}{261}
\GuideSection{Bibliography}{266}

\GuideChapter{\textbf{XXV. The Existence Theorem}, by M.~Demazure}{267}
\GuideSection{1. Statement of the theorem}{267}
\GuideSection{2. Existence theorem: construction of a group chunk}{269}
\GuideSection{3. Existence theorem: end of the proof}{274}
\GuideSection{4. Appendix}{276}
\GuideSection{Bibliography}{277}

\GuideChapter{\textbf{XXVI. Parabolic Subgroups of Reductive Groups}, by
  M.~Demazure}{279}
\GuideSection{1. Review. Levi subgroups}{279}
\GuideSection{2. Structure of the unipotent radical of a parabolic
  subgroup}{285}
\GuideSection{3. Scheme of parabolic subgroups of a reductive group}{289}
\GuideSection{4. Relative position of two parabolic subgroups}{296}
\GuideSection{5. Conjugacy theorem}{311}
\GuideSection{6. Parabolic subgroups and split tori}{320}
\GuideSection{7. Relative root datum}{322}
\endgroup

\clearpage
\begingroup
\small
\GuideSection{Bibliography}{330}
\endgroup

\clearpage
\thispagestyle{empty}
\null
